\section{Conclusion}
\label{ch:conclusion}

A 5\% reduction in carried-state persistence was the only tested RNN
perturbation to reproduce the complete descriptive majority pattern after
exploratory nomination and retest on distractor-trained circular networks.
The strongest present evidence is selective 2-back impairment across all ten
N-back checkpoints. Circular slowing, delay, and distractor contrasts had the
predicted average direction in four of five checkpoints, but their effects were
small and their seed-level intervals included zero. Training with distractors
removed the previous out-of-distribution confound without converting those
circular contrasts into robust findings.

Reduced persistence therefore remains a credible candidate for matched-cost
specificity and dynamical analysis, but the present result is an abstract
descriptive sufficiency finding against the native baseline---carried mainly by
the load dissociation---not evidence for a biological mechanism of psilocybin.

\begin{scaffold}
Rewrite after Sections~\ref{ch:results-specificity} and
\ref{ch:results-dynamics} are complete. Target story:
\begin{quote}
We built and validated working-memory RNNs, established how generic variability
disrupts them, derived candidate dynamical perturbations from psychedelic
theory and human task evidence, and tested which, if any, reproduces the
selective behavioural and representational signature better than nonspecific
noise.
\end{quote}
\end{scaffold}
